\section*{Chapter \ref{ch:b1:rate-laws} --- Chemical Kinetics: Rate Laws}

\begin{solution}{exo:b1:rate-laws:1}
$v = \frac14 \times 2.0 \times 10^{-4} = \qty{5.0e-5}{mol.L^{-1}.s^{-1}}$;
dioxygen forms at the same rate, \qty{5.0e-5}{}; \ce{N2O5} disappears at
$2v = \qty{1.0e-4}{mol.L^{-1}.s^{-1}}$.
\end{solution}

\begin{solution}{exo:b1:rate-laws:2}
Order 0: \unit{mol.L^{-1}.s^{-1}}; order 1: \unit{s^{-1}}; order 2:
\unit{L.mol^{-1}.s^{-1}}; order 3: \unit{L^2.mol^{-2}.s^{-1}}.
\end{solution}

\begin{solution}{exo:b1:rate-laws:3}
$t_{1/2} = \ln2/k = 0.693/(3.0 \times 10^{-3}) = \qty{231}{s}$. After
\qty{600}{s}, $\eu^{-kt} = \eu^{-1.8} = 0.165$: 16.5\,\% left.
\end{solution}

\begin{solution}{exo:b1:rate-laws:4}
A \omterm{def:b1:rate-laws:half-life}{half-life} proportional to $1/a_0$ means order 2. If halving $a_0$
halves the \omterm{def:b1:rate-laws:half-life}{half-life}, $t_{1/2} \propto a_0$: order 0.
\end{solution}

\begin{solution}{exo:b1:rate-laws:5}
$[\ce{A}]$ is not linear in $t$ (losses 0.0259, 0.0192, 0.0142, 0.0106).
$\ln[\ce{A}] = -2.303$, $-2.602$, $-2.902$, $-3.202$, $-3.503$: it drops by
0.300 every \qty{10}{min}, a straight line. Order 1, $k =
\qty{0.0300}{min^{-1}}$.
\end{solution}

\begin{solution}{exo:b1:rate-laws:6}
Doubling $[\ce{A}]_0$ multiplies $v_0$ by $4 = 2^2$: order 2 in \ce{A}.
Doubling $[\ce{B}]_0$ multiplies it by 2: order 1 in \ce{B}. $v =
k[\ce{A}]^2[\ce{B}]$ with $k = 2.0 \times 10^{-6}/(0.010^2 \times 0.010) =
\qty{2.0}{L^2.mol^{-2}.s^{-1}}$.
\end{solution}

\begin{solution}{exo:b1:rate-laws:7}
\ce{B} is in hundred-fold excess: its concentration stays
\qty{0.50}{mol/L}, and $v = k'[\ce{A}]$ with $k' = k[\ce{B}]_0$. Hence
$k = 0.012/0.50 = \qty{0.024}{L.mol^{-1}.s^{-1}}$.
\end{solution}

\begin{solution}{exo:b1:rate-laws:8}
$E_a = R\ln5/(1/300 - 1/320) = 8.314 \times 1.609/(2.083 \times 10^{-4}) =
\qty{64.2}{kJ/mol}$. At \qty{310}{K}: $\ln k = \ln(1.0 \times 10^{-3}) +
\frac{E_a}{R}\big(\frac{1}{300} - \frac{1}{310}\big) = -6.908 + 0.831$,
$k = \qty{2.3e-3}{s^{-1}}$.
\end{solution}

\begin{solution}{exo:b1:rate-laws:9}
Here $a = 2$: $1/[\ce{A}] = 1/a_0 + 2kt$, $t_{1/2} = 1/(2k a_0) = 1/(2
\times 0.50 \times 0.020) = \qty{50}{s}$. At \qty{100}{s}: $1/[\ce{A}] = 50
+ 100 = \qty{150}{L/mol}$, $[\ce{A}] = \qty{6.7e-3}{mol/L}$.
\end{solution}

\begin{solution}{exo:b1:rate-laws:10}
$[\ce{A}] = a_0(1 - f) = a_0\eu^{-akt_f}$ gives $t_f = -\ln(1-f)/(ak)$, and
in half-lives $t_f/t_{1/2} = -\ln(1-f)/\ln2$: 1 for 50\,\%, 2 for 75\,\%,
3.32 for 90\,\%, 9.97 for 99.9\,\%: about ten half-lives for the reaction
to be ``finished''.
\end{solution}

\begin{solution}{exo:b1:rate-laws:11}
$m(t) = 10 - 0.40\,t$ (\unit{mg}, \unit{h}): order 0. \omterm{def:b1:rate-laws:half-life}{Half-life}
$10/(2 \times 0.40) = \qty{12.5}{h}$; empty after \qty{25}{h}. The dose
delivered per hour does not depend on what is left: a steady supply, which
is the purpose of a patch.
\end{solution}

\begin{solution}{exo:b1:rate-laws:12}
$E_a = R\ln2/(1/298.15 - 1/308.15) = 8.314 \times 0.693/(1.088 \times
10^{-4}) = \qty{52.9}{kJ/mol}$. Between 273.15 and \qty{283.15}{K}:
$\exp\big(\frac{52900}{8.314}(\frac{1}{273.15} - \frac{1}{283.15})\big) =
2.3$: the same \omterm{def:b1:rate-laws:activation-energy}{activation energy} gives a larger factor at lower
temperature.
\end{solution}

\begin{solution}{pb:b1:rate-laws:1}
\textbf{1.} Losses 8.6, 7.9, 7.2 and 6.5 points: not constant, so not
order 0.
\textbf{2.} $\ln$: 4.605, 4.515, 4.425, 4.335, 4.245; it drops by 0.090
every 10 days: a straight line, order 1.
\textbf{3.} $k_{60} = 0.090/10 = \qty{9.0e-3}{d^{-1}}$.
\textbf{4.} $t_{1/2} = \ln2/k_{60} = \qty{77}{d}$.
\textbf{5.} $100\,\eu^{-0.54} = 58.3\,\%$.
\textbf{6.} $t_{90} = \ln(10/9)/k_{60} = 0.1054/0.0090 = \qty{11.7}{d}$.
\textbf{7.} 90\,\% left: $\eu^{-kt_{90}} = 0.9$, so $t_{90} = \ln(10/9)/k$.
\textbf{8.} $t_{90}/t_{1/2} = \ln(10/9)/\ln2 = 0.152$.
\textbf{9.} For order 1 the fraction left depends on $kt$ only, not on the
initial amount.
\textbf{10.} At \qty{25}{\celsius} the degradation takes years; heating
accelerates it (Arrhenius) so that it can be measured in weeks.
\textbf{11.} For order 2 the \omterm{def:b1:rate-laws:half-life}{half-life} is $1/(ak\,a_0)$: a tablet twice as
strong would degrade twice as fast in relative terms and keep less long.
\textbf{12.} $1/T$ (\unit{K^{-1}}): \num{3.1934e-3}, \num{3.0945e-3},
\num{3.0017e-3}; $\ln k$: $-6.669$, $-5.661$, $-4.711$.
\textbf{13.} Slope $= (-4.711 + 6.669)/(3.0017 \times 10^{-3} - 3.1934 \times
10^{-3}) = \qty{-1.021e4}{K}$.
\textbf{14.} $E_a = 8.314 \times 1.021 \times 10^4 = \qty{85}{kJ/mol}$.
\textbf{15.} Predicted $\ln k_{50} = -4.711 - 1.021 \times 10^4 \times (3.0945 -
3.0017) \times 10^{-3} = -5.658$, measured $-5.661$: on the line.
\textbf{16.} $A = k_{60}\,\eu^{E_a/RT} = 0.0090\,\eu^{10215/333.15} =
\qty{1.9e11}{d^{-1}}$.
\textbf{17.} $\exp\big(1.021 \times 10^4(\frac{1}{298.15} -
\frac{1}{308.15})\big) = 3.0$: more than the factor 2 of the rule, because
$E_a$ is larger than \qty{53}{kJ/mol}.
\textbf{18.} $k_{30} = k_{60}\exp\big(-1.021 \times 10^4(\frac{1}{303.15} -
\frac{1}{333.15})\big) = \qty{4.3e-4}{d^{-1}}$.
\textbf{19.} $t_{90} = 0.1054/4.33 \times 10^{-4} = \qty{243}{d}$, 8.0
months.
\textbf{20.} $1 - \eu^{-3 \times 0.0090} = 2.7\,\%$.
\textbf{21.} The shelf life falls quickly with temperature (8 months at
\qty{30}{\celsius} against 14 at \qty{25}{\celsius}): the expiry date is
valid only below \qty{25}{\celsius}.
\textbf{22.} $k_{25} = k_{60}\exp\big(-1.021 \times 10^4(\frac{1}{298.15} -
\frac{1}{333.15})\big) = \qty{2.46e-4}{d^{-1}}$.
\textbf{23.} $t_{90} = 0.1054/2.46 \times 10^{-4} = \qty{428}{d}$, that
is \textbf{14 months}: the date printed on the box, obtained from six
weeks of measurements.
\end{solution}
